\section{算法实现与计算流程}
\subsection{公共数据结构与计算接口}
\paragraph{对象化数据组织}
为保证四问口径一致，程序不直接在各问题中重复读取工作簿，而是先构造节点、货箱、机型、实体资源、航段和通信端点六类对象。节点对象保存经纬度、地面海拔与作业高度；货箱对象保存目的地、质量、体积、类别、优先系数及适用时限；机型对象保存容量、速度、航程、能量与作业时间参数；实体资源对象保存类型、编号和下一可用时刻。

航段计算由起终点、水平距离、穿越DEM像元、最高地形、巡航海拔、爬升和下降高度组成。当前实现通过同一物理函数统一计算这些量，但没有把节点对几何结果持久缓存；不同机型和载荷会重新调用该函数，以保证时间与能耗口径一致。通信端点在此基础上增加天线高度和随时间变化的位置，使运输与通信使用相同的地形和坐标口径。

\paragraph{统一路线评价函数}
任意运输候选均交给同一个路线评价函数。输入为机型、货箱集合、访问顺序和准备开始时刻，函数按以下顺序计算：
\begin{enumerate}
  \item 汇总起飞质量与体积，检查额定容量；
  \item 根据准备和逐箱装载时间得到起飞时刻；
  \item 逐航段读取地形、计算飞行时间与当前剩余载荷能耗；
  \item 到站后按硬截止、优先系数和编号确定交接次序；
  \item 记录每箱交付完成时刻并扣除已交付质量；
  \item 计算末站返航、总能耗、返航SOC和电池充电结束时刻；
  \item 返回可行标志、全部事件和目标指标。
\end{enumerate}
问题一把路线限制为单点往返，问题二允许多点序列，问题三在相同运输事件上叠加通信检查。因此，不同问题不会因为使用不同路线函数而产生能耗或时间口径差异。

\subsection{问题一求解器实现}
\paragraph{安全载荷与候选箱组}
对每个“服务区—机型”组合先执行空载和满载边界检查，再对内部边界进行60次二分。随后将该服务区的$n_i$个货箱编码为$n_i$位掩码，预先计算各掩码的质量、体积和货箱数。一个掩码只有在至少一种机型下同时通过容量和能量检查时才进入候选表。

候选表按掩码索引，保存所有可执行机型及对应能耗和作业时长。实现中不另设显式支配删除；同一覆盖状态只保留字典序代价更小的路径，因此较差候选在动态规划更新时自然淘汰，同时保留不同覆盖集合之间可能产生的组合优势。

\paragraph{状态递推与结果恢复}
动态规划状态为已经覆盖的货箱掩码$M$。从剩余货箱的最低有效位选择锚定货箱，只枚举包含该位且与$M$不相交的候选子集$S$，以避免同一批次集合因顺序不同被重复搜索。每个状态直接保存最优字典序代价及截至该状态的完整候选路径；全集状态中的路径就是最终批次。

得到完整路径后，程序汇总全部货箱编号并执行唯一覆盖审计；同时，候选进入动态规划前已经由统一物理函数完成质量、体积和返航能量检查。若某箱未被覆盖、重复覆盖或不存在全集可行状态，则整套方案判为无效。该检查用于验证状态压缩结果与原始需求是否一致。

\subsection{问题二调度器实现}
\paragraph{两套确定性任务构造}
问题二构造并比较两套方案。第一套为硬箱专送基线：先形成硬时限任务，再对剩余普通货箱按问题一方法组批，并把能够节能的两个异区单点任务贪心合并为双区串访。第二套为同目的地混装方案：在保持基线紧急任务实体机指派的条件下，枚举同目的地软箱补入子集，再对剩余货箱重复精确组批和异区串访合并。双区路线同时评价两种访问顺序，只有完整路线函数返回容量和能量可行时才进入候选集。

两套候选均完整执行实体资源排程，并以$(N_{late}^{hard},WTD,C_{max},E_T,N_T)$字典序选择结果。当前程序没有实施通用的箱组拆分、同区合并或反复局部邻域迭代，因此所得方案是两种确定性构造中的较优者，而不是局部搜索收敛解或全局最优解。

\paragraph{实体机—电池事件排程}
对每类机型维护实体机最早可用时刻和同型电池满电时刻。排入一个任务时，枚举全部同型实体机—电池组合，以二者可用时刻的最大值作为最早准备开始时刻，再调用路线评价函数。各组合的硬迟到箱数、最大硬迟到、软时限代价、交付时刻、能耗和资源恢复时刻共同进入评分；程序选择评分最小的组合，并更新实体机返航时刻与电池充满时刻。硬迟到组合不会在生成时被直接删除，但会被零硬迟到组合严格支配；最终方案另行核验硬迟到为0。

全部任务排入后，把每个开始和结束事件写入按时间排序的事件流。扫描实体机区间和电池区间时，结束事件在相同时刻先于开始事件处理，保证半开区间约定一致。最后从逐箱交付表计算硬违约、软超期和$WTD$，从任务表计算完工时间、能耗和架次数，避免在搜索过程中维护的增量目标与最终明细不一致。

\subsection{问题三连续通信实现}
\paragraph{离线选点与固定波次输入}
运输路线被展开为准备后的爬升、水平巡航、下降、站内交接和返航区间。飞行区间位置按时间线性变化，交接区间位置固定。六个中继点由独立选点脚本在5~s黑障样本上离线搜索并筛选，随后将点位和W1--W7七个波次固化到主求解器；主程序要求输入恰为所选问题二的22个运输任务。这个绑定关系保证结果可复现，但也意味着当前实现不是每次运行时在线生成覆盖矩阵并贪心选点。

覆盖矩阵只服务于候选选择。选中点形成中继任务后，程序从每个运输时段读取实际中继编号，计算直连余量或两条中继链路余量的较小值，并生成通信分段。相邻且通信方式、实际中继和可用性依据相同的时段合并，最终得到68个可解释分段。

\paragraph{连续区间递归认证}
连续认证函数的输入为时间区间、运输轨迹段和候选通信方式。它先计算区间上的保守余量下界；下界非负则返回“整段通过”。若无法证明且区间长度大于0.1~s，则在中点二分并分别认证；任一子区间失败，父区间即失败。达到最小长度仍无法证明时按未认证处理，不用端点或中点样本替代证明。

递归返回的不仅是通过标志，还包括最小下界余量及其对应架次和区间。对最终方案再次按0.5~s步长展开全部运输轨迹，独立统计可用点与中断点。连续证书和网格检查分别写入结果表，二者均通过才把运输架次标记为通信可行。

\subsection{问题四分区枚举实现}
\paragraph{原子块与规范化分组}
先以服务区为顶点扫描全部运输路线：同一多点路线访问的服务区之间加入共访边，通过并查集或深度优先搜索得到连通分量。每个连通分量成为一个原子块，并保存其服务区集合、关联运输任务及由通信分段反推的中继任务。

枚举时按固定原子块顺序赋组，并把第一个原子块固定在第1组；其余原子块可取任一组标签，最后要求$K$个组均非空。由此实际遍历1023个两组候选，以及57002个三组带标签候选；后三者中第2、3组标签互换形成等价对，对应28501个无标签等价类。

\paragraph{资源扫描与方案排序}
对每个组收集分类型任务区间。运输机和电池按A/B/C型分别扫描，中继机与能源组件单独扫描；电池区间包含充电，中继机区间包含周转。将开始事件记为$+1$、结束事件记为$-1$，按“时刻、结束优先”的顺序累计，最大累计值即该类型最小配置。

各组资源向量汇总后，程序按“总库存缺口、总配置数、工作量变异系数、完成时刻极差、规范化分组”进行字典序排序；相对集中执行的冗余在选定方案后逐资源计算并报告，不进入当前排序键。程序保留最优方案的完整分类型向量与每组服务区，避免仅输出总件数而无法判断究竟缺少哪类资源。

\subsection{关键实现参数与输出接口}
表\ref{tab:implementation-parameters}列出影响数值判断和搜索边界的关键参数。参数要么直接来自题目附件，要么用于控制数值精度；算法没有通过调整这些参数改变题设物理条件。

\begin{table}[htbp]\centering
\caption{关键实现参数及其作用}\label{tab:implementation-parameters}
\small
\begin{tabularx}{\textwidth}{>{\raggedright\arraybackslash}p{3.0cm}YY}
\toprule
参数或约定 & 取值与来源 & 计算作用\\
\midrule
DEM基础分辨率 & 30~m，题目附件 & 地形净空与视线遮挡的空间基础\\
航段密集采样间隔 & 15~m，并补充像元边界交点 & 辅助枚举经过像元，不替代像元遍历\\
安全载荷二分次数 & 60次，数值精度设置 & 使质量区间远小于输出三位小数精度\\
资源区间约定 & 左闭右开、同刻结束优先 & 避免相邻任务在边界被误判重叠\\
连续认证最小区间 & 0.1~s，保守认证设置 & 未能证明的短区间按不可用处理\\
通信独立复核步长 & 0.5~s & 对连续证书结果进行高密度抽查\\
目标比较方式 & 硬可行优先的字典序 & 避免不同量纲指标由任意权重相加\\
\bottomrule
\end{tabularx}\end{table}

输出分为逐箱、逐运输架次、逐中继架次、逐通信分段和逐分区五个层次。逐箱记录用于核对覆盖和时限；架次记录用于回算能量和资源；通信分段用于证明连续保障并确定问题四的实际中继归属；分区记录保留分类型配置、缺口、冗余与均衡指标。分层输出使摘要和汇总表中的指标都能够由更细一级记录重新计算。
